\documentclass[11pt]{article}
\input{../../shared/project_style.tex}
\rhead{Basics 17 Textbook Draft}

\begin{document}
\DocTitle{Basics 17: What Is a Plasma?}{III - Elementary plasma physics}{Ionization, quasineutrality, collective behavior, screening, coupling, and plasma regimes}

\begin{overviewbox}
\textbf{Textbook draft, version 1.0.} A plasma is not merely a hot gas containing ions. Its defining behavior is collective: many charged particles reorganize electric and magnetic fields over characteristic lengths and times. This chapter develops practical criteria for deciding when an ionized medium should be modeled as a plasma. This chapter is designed for a reader with no prior plasma-physics training. It distinguishes exact identities from physical approximations and connects the central equations to later simulation modules.
\end{overviewbox}

\ProjectTOC

\section{Learning objectives}
After completing this note, the reader should be able to:
\begin{itemize}[leftmargin=1.6em]
\item Define ionization fraction, quasineutrality, and collective behavior.
\item Use Debye length and plasma parameter as plasma criteria.
\item Distinguish weakly and strongly coupled plasmas.
\item Explain why local charge separation can coexist with macroscopic quasineutrality.
\item Classify common laboratory, space, and fusion plasmas.
\item Identify when fluid, kinetic, or neutral-gas descriptions may be appropriate.
\end{itemize}

\section{Prerequisites and reading strategy}
\begin{itemize}[leftmargin=1.6em]
\item Basics 1, 10-13.
\end{itemize}
On a first reading, emphasize physical meaning and dimensional checks. On a second reading, reproduce the derivations. Attempt the computational laboratory only after the limiting cases are understood.

\section{Ionized matter is necessary but not sufficient}

A gas containing electrons and positive ions has ionization fraction
\begin{equation}
\chi_i=\frac{n_i}{n_i+n_n},
\end{equation}
where $n_n$ is neutral density. A medium can exhibit plasma behavior even when $\chi_i\ll1$ if charged particles respond collectively before neutral collisions erase that response. Conversely, a sparse collection of noninteracting charged particles is not automatically a plasma.


\section{Collective behavior}

Each charge interacts through long-range electromagnetic fields. In a plasma, the relevant response is not a sequence of isolated binary encounters but the coordinated rearrangement of many particles. Density disturbances can create plasma oscillations; currents can generate magnetic fields; pressure and magnetic tension can launch waves. Collective dynamics is meaningful when many particles occupy a shielding volume and system dimensions exceed the shielding length.


\section{Quasineutrality and local charge separation}

For singly charged ions,
\begin{equation}
n_e\simeq n_i,
\end{equation}
so the net charge density
\begin{equation}
\rho_q=e(n_i-n_e)
\end{equation}
is small compared with either species charge density. Quasineutrality is an asymptotic statement on scales much larger than the Debye length. It does not require exact equality point by point. Sheaths, double layers, waves, and small charge imbalances can generate important electric fields.


\section{Debye length and system-size criterion}

For electrons of temperature $T_e$ and density $n_e$, the Debye length is
\begin{equation}
\lambda_D=\sqrt{\frac{\epsilon_0 k_BT_e}{n_ee^2}}.
\end{equation}
A common plasma criterion is
\begin{equation}
L\gg\lambda_D,
\end{equation}
so that an interior region can remain quasineutral while charge imbalances are screened over smaller scales. When $L$ is comparable to $\lambda_D$, boundaries and electrostatic kinetics dominate.


\section{Particles in a Debye sphere}

The number of particles in a Debye sphere is
\begin{equation}
N_D=\frac{4\pi}{3}n_e\lambda_D^3.
\end{equation}
The criterion $N_D\gg1$ means that screening is a collective average over many particles and fluctuations are relatively small. It is closely related to weak coupling. Some texts define the plasma parameter using $n\lambda_D^3$ without the geometric factor; conventions should be stated.


\section{Weak and strong coupling}

Compare a characteristic Coulomb potential energy with thermal energy:
\begin{equation}
\Gamma=\frac{e^2}{4\pi\epsilon_0 a k_BT},\qquad
a=\left(\frac{3}{4\pi n}\right)^{1/3}.
\end{equation}
For $\Gamma\ll1$, thermal motion dominates pair interaction energy and conventional weakly coupled plasma theory applies. For $\Gamma\gtrsim1$, correlations become strong; dusty plasmas, dense matter, and some ultracold plasmas require different methods.


\section{Magnetization and collisionality}

A species is magnetized when it completes many gyrotations before collisions significantly scatter it and when its Larmor radius is small relative to system gradients:
\begin{equation}
|\Omega_s|\gg\nu_s,\qquad \rho_s\ll L.
\end{equation}
Electrons and ions can occupy different regimes simultaneously. Plasma classification therefore requires species, scale, and timescale labels rather than one global adjective.


\section{Examples and modeling choices}

Fusion plasmas are hot, magnetized, and often weakly collisional. The solar wind is dilute and collisionless but strongly collective. Fluorescent discharges can be partially ionized and collisional. Stellar interiors may be dense and collisional. A suitable model is selected by comparing $\lambda_D$, $\rho_s$, mean free paths, wave periods, and device dimensions.


\section{Computational laboratory}
Create a regime calculator that accepts density, temperature, magnetic field, system size, and collision frequency. Compute $\lambda_D$, $N_D$, $\Gamma$, gyrofrequencies, Larmor radii, and magnetization ratios. Apply it to a fusion core, glow discharge, solar wind, and ionosphere case.

\section{Summary}
\begin{itemize}[leftmargin=1.6em]
\item Plasma behavior is defined by collective electromagnetic response, not ionization alone.
\item Quasineutrality is valid on scales large compared with the Debye length.
\item Many particles per Debye sphere support a smooth collective description.
\item The coupling parameter separates weakly and strongly correlated regimes.
\item Model selection must be made for specific species, lengths, and times.
\end{itemize}

\SymbolsAcronymsSection
\begin{symtable}
$\chi_i$ & ionization fraction & Fraction of particles in the ionized component. \\
$\lambda_D$ & Debye length & Electrostatic screening scale. \\
$N_D$ & Debye-sphere population & Particles inside a shielding volume. \\
$\Gamma$ & coupling parameter & Potential-to-thermal energy ratio. \\
$a$ & Wigner-Seitz radius & Characteristic interparticle spacing. \\
$\nu_s$ & collision frequency & Rate of species-$s$ scattering. \\
\end{symtable}

\section{Exercises}
\begin{enumerate}[leftmargin=1.8em]
\item Explain why an ionized gas with $L\ll\lambda_D$ may fail the usual plasma criterion.
\item Compute $N_D$ for a specified electron density and temperature.
\item Show how $\lambda_D$ changes when density is increased by a factor of four.
\item Distinguish quasineutrality from zero electric field.
\item Discuss why electrons may be magnetized while ions are not.
\end{enumerate}

\section{References}
\begin{thebibliography}{99}
\bibitem{ref1} D. J. Griffiths, \emph{Introduction to Electrodynamics}, 4th ed., Cambridge University Press (2017).
\bibitem{ref2} J. D. Jackson, \emph{Classical Electrodynamics}, 3rd ed., Wiley (1998).
\bibitem{ref3} F. F. Chen, \emph{Introduction to Plasma Physics and Controlled Fusion}, 3rd ed., Springer (2016).
\bibitem{ref4} R. J. Goldston and P. H. Rutherford, \emph{Introduction to Plasma Physics}, CRC Press (1995).
\bibitem{ref5} P. M. Bellan, \emph{Fundamentals of Plasma Physics}, Cambridge University Press (2006).
\end{thebibliography}

\end{document}
